\clearpage
\VerdierPhysicalPageAnchor{0191}%
\begingroup\small\setlength{\parskip}{0pt}\setlength{\emergencystretch}{1.5em}
\setlength{\abovedisplayskip}{4pt plus 2pt minus 2pt}%
\setlength{\belowdisplayskip}{4pt plus 2pt minus 2pt}%

\VerdierSemanticAnchor{ch3d-2-2-2}
\noindent\textbf{引理 2.2.2.}\quad
\textit{设
$0\longrightarrow X\longrightarrow Y\longrightarrow Z\longrightarrow0$
是 $\mathcal A$ 中的正合列。}

\noindent\textit{a) 若 $X$ 与 $Y$ 是 $M^{\wedge}$ 的元素，则存在满态射
$u\colon L\longrightarrow Z$，其中 $L\in M$。}

\noindent\textit{b) 若 $X$ 与 $Z$ 是 $M^{\wedge}$ 的元素，则存在满态射
$u\colon L\longrightarrow Y$，其中 $L\in M$。}

\noindent\textit{c) 若 $Y$ 与 $Z$ 是 $M^{\wedge}$ 的元素，则存在满态射
$u\colon L\longrightarrow X$，其中 $L\in M$。}

\VerdierSemanticAnchor{ch3d-2-2-3}
\medskip\noindent\textbf{引理 2.2.3.}\quad
\textit{设
$0\longrightarrow X\longrightarrow Y\longrightarrow Z\longrightarrow0$
是 $\mathcal A$ 中的正合列。}

\noindent\textit{a) 若 $Y$ 与 $Z$ 是 $M^{\wedge}$ 的元素，则对每个满态射
$u\colon L\longrightarrow X$，其中 $L\in M$，$u$ 的核是
$M^{\wedge}$ 的两个元素之间某个满态射的核。}

\noindent\textit{b) 若 $X$ 与 $Y$ 是 $M^{\wedge}$ 的元素，则对每个满态射
$u\colon L\longrightarrow Z$，其中 $L\in M$，$u$ 的核是
$M^{\wedge}$ 的两个元素之间某个单态射的余核。}

\noindent\textit{c) 若 $X$ 与 $Z$ 是 $M^{\wedge}$ 的元素，则对每个满态射
$u\colon L\longrightarrow Y$，其中 $L\in M$，$u$ 的核是
$M^{\wedge}$ 的两个元素之间某个满态射的核。}

\VerdierSemanticAnchor{ch3d-2-2-2-3-proof}
这两个引理的证明留给读者。请注意，引理 (2.2.2) 与 (2.2.3) 的性质
(a) 和 (c) 不需要性质 $(E_2)$。

\VerdierSemanticAnchor{ch3d-2-2-4}
\medskip\noindent\textbf{命题 2.2.4.}\quad
\textit{a) 集合 $M^{\wedge}$ 具有性质 $(E_6)$ 与 $(E_2)$。集合
$M$ 包含于 $M^{\wedge}$。$M^{\wedge}$ 的每个对象都是 $M$ 中某个对象的商。}

\noindent\textit{b) 设对象集合 $M'$ 包含 $M$，且 $M'$ 的每个对象都是
$M$ 中某个对象的商，并且具有性质 $(E_3)$；则该集合包含于 $M^{\wedge}$。}

\noindent\textit{c) 特别地，
$(M^{\wedge})^{\wedge}=M^{\wedge}$。}

\VerdierSemanticAnchor{ch3d-2-2-4-proof}
性质 a) 由前两个引理推出。性质 b) 与 c) 是直接的。

\VerdierSemanticAnchor{ch3d-2-2-5}
\medskip\noindent\textbf{2.2.5.}\quad
设 $M\subset M'$ 是 $\mathcal A$ 的两个对象集合，其中 $M$ 具有性质
$(E_4)$ 与 $(E_5)$，$M'$ 具有性质 $(E_4)$，而且 $M'$ 的每个对象都是
$M$ 中某个对象的商。对于对象 $X$（属于 $M'$），称其 $M$-维数小于或等于
$n$，并记作 $M\text{-}\dim(X)\leq n$，如果存在正合列

\VerdierSemanticAnchor{ch3d-2-2-5-exact-sequence}
\[
0\longrightarrow L_n\longrightarrow L_{n-1}\longrightarrow\cdots
\longrightarrow L_0\longrightarrow X\longrightarrow0,
\qquad L_i\in M,\quad0\leq i\leq n.
\]
\endgroup
